\documentclass[10pt]{beamer}
\usepackage{csquotes}% Recommended

\usepackage{graphicx}
\usepackage{xcolor}
\usepackage{amsfonts}
\usepackage{amsmath}
\usepackage{amssymb}
\usepackage[skip=10pt plus1pt] {parskip}
\usepackage{multicol}
\usepackage{multirow}

%Information to be included in the title page:
\title{Probability and Its Associated Calculus}
\author{UChicago Political Science Math Camp}
\institute{Ryan Gibbons \& Luke Cain}
\date{2025}

\begin{document}

\frame{\titlepage}


\begin{frame}
\begin{center}
    \includegraphics[width = 6cm]{bayesian (1).jpg}
\end{center}    
\end{frame}

\begin{frame}{Probability}

Probability is the \textbf{likelihood of an event occurring} \\

\vspace{1cm}

We denote the probability of an event $A$ as a real number between 0 and 1:

\[\Pr[A] \in [0,1]\]

\end{frame}


\begin{frame}{Discrete vs. Continuous Probability}
    There are two ``categories" in studying probability:
    \begin{enumerate}
        \item \textbf{Discrete} probability, where there are $n \in \mathbb{N}$ possible outcomes or events (``What number will this dice roll?")
        \item \textbf{Continuous} probability, where there are $n \to \infty$ possible outcomes or events (``How far will I throw this dice?")
    \end{enumerate}
    This refers to the \textit{measurement}, not the \textit{probability} itself
\end{frame}

\begin{frame}{Axioms of Probability}
    There are three basic axioms when working with probability:
    \begin{enumerate}
        \item \[\forall\; a \in A, \Pr[a] \in \mathbb{R}_+ \implies \Pr[a] \geq 0\] \pause

        \item \[\forall A,\; \sum_{a \in A} \Pr[a] = 1\] \pause

        \item \[a_i \cap a_{-i} = \emptyset \;\forall \;a_i, a_{-i} \iff \Pr\bigg[ \bigcup_{i=1}^k a_i\bigg] = \sum_{i=1}^k \Pr[a_i]\] \pauseG
    \end{enumerate}
\end{frame}

\begin{frame}{Applying Axioms}
    Show these axioms for each of the following discrete probabilities:

    \begin{enumerate}
        \item A fair two-sided coin
        \item An unfair six-sided dice
        \item Tomorrow's temperature (in whole degrees, ok to place bounds)
    \end{enumerate}
\end{frame}

\begin{frame}{Basic Probability}
    Let there be be an event with equally likely outcomes $A_{i \in n} \in \{S\}$. Then, the probability that a draw yields an outcome meeting some condition $ Q$ is:

    \[\Pr[a_i \rightarrow Q] = \frac{|a_i \rightarrow Q|}{|S|}\]

    \textit{Count of ``target" outcomes over count of total}
\end{frame}

\begin{frame}{Rolling a Dice}
    Suppose we roll a (fair) 100-sided dice to recover $a \in \{x \in \mathbb{N} : 1 \leq x \leq 100\}$. What is the probability of each outcome?

    \begin{enumerate}
        \item $a = 15$
        \item $a \leq 50$
        \item $a < 50$
        \item $a$ is even
        \item $\frac{a}{3} \in \mathbb{N}$
        \item $\frac{a}{7} \in \mathbb{Q}$
    \end{enumerate}
\end{frame}

\begin{frame}{Combinatorics}
    To determine $|S|$, we often rely on combinatorial mathematics, and use the formula:

    \[C = \begin{pmatrix}
        n \\
        a
    \end{pmatrix}= \frac{n!}{a!(n-a)!} \]

    Where
    \begin{itemize}
        \item $C$ = number of unique combinations (outcomes)
        \item $n$ = number of objects in a set
        \item $a$ = number of objects chosen from the set
    \end{itemize}
\end{frame}

\begin{frame}{Relaxed Assumptions}
    Of course, not all events are ``fairly" distributed:

    \begin{itemize}
        \item A cat landing on its feet
        \item An authoritarian election (mechanism design)
        \item Getting pulled over while speeding
    \end{itemize}

    Often, we can observe the outcomes, but must estimate its true probability
\end{frame}


\begin{frame}{Independence}
    Two events are said to be \textbf{independent} or \textbf{disjoint} if the probability of one is not tied to the probability of the other:

    \begin{itemize}
        \item Metra running late and the White Sox score
        \item Math background and field seminar grades
    \end{itemize}
\pause
    Two events are \textbf{dependent} if the expected outcome of one \textit{is} affected by the outcome of the other:

    \begin{itemize}
        \item Delays on the 171 and delays on the 172
        \item Numbered street and apartment rent
    \end{itemize}
\end{frame}

\begin{frame}{Conditional Probability}
    Many events are contingent on prerequisite outcomes. The \textbf{conditional probability} of an outcome $A$ given an event $B$ can be expressed as:

    \[\Pr[A \; | \; B] = \frac{\Pr[ A \cap B]}{\Pr[B]}\]
\end{frame}

\begin{frame}{A Note on Dependency}
    A dependent relationship can exist in different mechanisms:

    \begin{itemize}
        \item \textbf{Correlational}: The outcome of one event does not directly impact another, but signals underlying conditions
        \item \textbf{Causal}: The outcome of one event is causally related to the outcome of another
    \end{itemize}

    \textit{For now, our quantitative mechanisms are the same for each.}
\end{frame}

\begin{frame}{Replacement}
    One common example of independence versus dependence is whether or not an outcome is ``replaced" after a random draw:

    \begin{itemize}
        \item A card is drawn from a deck, then shuffled back in $\rightarrow$ \textbf{independent}
        \item A card is drawn from a deck, then discarded $\rightarrow$ \textbf{dependent}
    \end{itemize}
\vspace{.5cm}
    Why might this be?
\end{frame}

\begin{frame}{I.I.D.}
    Independent and identically distributed:

\begin{itemize}
    \item \underline{Independent}: The expected outcome of the $n$th event is not impacted by any $n-k$th outcome.
    \item \underline{Identically Distributed}: Every draw is identical, and occurs with full replacement
\end{itemize}
\end{frame}




\begin{frame}{Joint Probability}
    For two independent and disjoint events $A$ and $B$, their \textbf{joint probabilities} are as follows:

    \[\Pr[A \cap B] = \Pr[A] \cdot \Pr[B]\]

    \[\Pr[A \cup B] = \Pr[A] + \Pr[B]\]

    How might this change for non-disjoint events? \pause

    \[\implies \Pr[A \cup B] = \Pr[A \cup \neg B] + \Pr[\neg A \cup B] + \Pr[A \cap B]\]
    
\end{frame}

\begin{frame}{Joint Probability Example}
We want to evaluate each quadrant for $\Pr[i \cap j]$:
    \begin{table}[H]
\centering
\renewcommand{\arraystretch}{1.5} % Adjust row height
\setlength{\tabcolsep}{12pt} % Adjust column width
\begin{tabular}{cccc}
                    &                        & \multicolumn{2}{c}{$j$}                              \\
                    &                        & Sunny                        & Cloudy                        \\ \cline{3-4} 
\multirow{1}{*}{$i$} & \multicolumn{1}{c|}{Rainy} & \multicolumn{1}{c|}{} & \multicolumn{1}{c|}{} \\ \cline{3-4} 
                    & \multicolumn{1}{c|}{Dry} & \multicolumn{1}{c|}{} & \multicolumn{1}{c|}{} \\ \cline{3-4} 
\end{tabular}
\end{table}
\end{frame}


\begin{frame}{Expanding Conditional Probability}
    
\end{frame}

\begin{frame}{Suppose...}
    \begin{itemize}
        \item 
    \end{itemize}
\end{frame}


\begin{frame}{Bayesian Statistics}
    Suppose we know the probability of $A$ given $B$, but we want to know the probability of $B$ given $A$. We can use \textbf{Bayes' Rule}:

    \[\Pr[A \; | \; B] = \frac{\Pr[B \; | \; A] \cdot \Pr[A]}{\Pr[B]}\]
\end{frame}

\begin{frame}{Random Variables}
    A random variable is a function $X$ such that there exists a real number for every $s \in S$:

    \[X: S \rightarrow \mathbb{R}\]

    There is then some non-deterministic probability $\Pr[X(s) = x] \in [0, 1]$ for every event $s$.
\end{frame}

\begin{frame}{PMFs}
    For a discrete random variable $X$, we can model a \textbf{probability mass function}:

    \[\text{PMF} \equiv f(x) = \Pr[X = x]\]

    This models the likelihood of any given outcome of $x \in S$. \\

    \vspace{.5cm}
    What do our probability axioms imply about PMFs?
\end{frame}

\begin{frame}{PMF examples}
    How might we construct the PMFs of the following?

    \begin{itemize}
        \item A fair, six-sided dice
        \item An unfair six-sided dice (twice as likely to land on 6)
        \item A nation's global ranking in GDP (assume $n = 200$ for simplicity
    \end{itemize}
\end{frame}

\begin{frame}{CDFs}
    We often want to represent the probability that a random variable is larger or smaller than some threshold. For this, we can model a \textbf{cumulative distribution function}:

    \[\text{CDF} \equiv F(x) = \Pr[X \leq x] = \sum_{i \leq x} \Pr[i]\]

    What do our probability axioms imply about discrete CDFs? How might we solve $\Pr[X \geq x]?$
\end{frame}

\begin{frame}{CDF examples}
    How might we construct CDFs of the following?

    \begin{itemize}
        \item A fair dice rolling at least a 3
        \item A nation being in the top $n$th percentile globally of GDP (assume $N = 200$ for simplicity)
        \item The ``amount" of persuasion voters need to choose candidate $A$ in a two-race, majoritarian election
        
    \end{itemize}
\end{frame}

\begin{frame}{Uniform Distribution}
    The \textbf{uniform distribution} models every outcome as equally likely:

    \[f(x) = \Pr[X = x] = \frac{1}{|X|} \]

    Note that $f(x)$ is a \textit{constant function} -- what does this imply?
\end{frame}

\begin{frame}{Another Important Discrete Distribution}
    Suppose you flip a fair coin twice. What is the probability of recovering $H$ at least once? \\

    \vspace{.5cm}

    Suppose you roll a fair dice six times. What is the probability of recovering $6$ at least once? \\

\end{frame}

\begin{frame}{A Closing Derivation}
    Let each draw of an event consist of $n$ equally likely outcomes across $N$ alternatives. What is the probability of drawing some outcome $n \equiv Q$ over the period of draws $1, 2, ...,t$ for a condition $Q$?
\end{frame}

\begin{frame}{The Binomial Distribution}
To model the likelihood of an event of probability $p$ occurring within $n$ i.i.d. trials, we can use the \textbf{binomial distribution}:

    \[B(n,p) = \begin{pmatrix}
        n \\ k
    \end{pmatrix} p^k q^{n-k}\]

    Where $\begin{pmatrix}
        n \\k
    \end{pmatrix}$ is the \textbf{binomial coefficient} for $k$ outcomes over $n$ trials
\end{frame}

\begin{frame}{Expectation (A Minor Recap)}
    The \textbf{expected value} of a random variable is the sum of its possible outcome values weighted by the probability that value occurs:

    \[E[X] = \sum_x \Pr[x] \cdot x\]

    With discrete RVs, this summation is simple and well-defined, even if not always meaningful
\end{frame}

\begin{frame}{Expectation Examples}
    Calculate the expected value of the following:
    \begin{enumerate}
        \item A fair six-sided dice roll
        \item The sum of two fair six-sided dice rolls
        \item The \textbf{maximal value} of two fair dice rolls
    \end{enumerate}
\end{frame}

\begin{frame}{Extending Expectation}
    Let's evaluate this. Google ``Roll Two Dice", and record 10 iterated outcomes each.\\ \pause
\vspace{1cm}
What do we start to see when plotting these outcomes?
    
\end{frame}

\begin{frame}{And More So, Large $N$}
    What if we could simulate a more sophisticated result -- the maximal value of three dice rolls?
\end{frame}




\end{document}